% Incorporación aditiva. No modifica las fuentes ni el PDF de 211 páginas.
% Procedencia, dependencias y alcance en preservation_map.json.
\section{Représentation grassmannienne des lecteurs et formes à valeurs dans la droite d’action}
\label{sec:compat-lectores-formas}

Les réalisations positives conservent une information qui peut être réutilisée
par les lecteurs physiques. Cette conservation possède une formulation
précise : une représentation marquée du registre permet de reconstruire ses
coordonnées et de transporter, par composition, chaque fonction qui les utilise.
La géométrie représente un résultat de la construction ; son choix de
coordonnées conserve la provenance de ce résultat.

Le domaine de départ est formé des états produits par APP, TRIT et TPK dans les sections précédentes. APP conserve ses feuillets additif et multiplicatif au moyen du résidu et du quotient ; TRIT sélectionne le régime local avec son orientation ; la composition \(U_t=\operatorname{Upd}_t\circ\operatorname{Tra}_t\circ
\operatorname{Sel}_t\) met à jour la position, le registre et la mémoire. Le prolongement
\[
 w_6\longrightarrow w_{12}\longrightarrow w_{18}\longrightarrow
 w_{24}\longrightarrow w_{30}\longrightarrow R_{36}\longrightarrow G_9
\]
conserve le préfixe et distingue retour de phase et retour de l'état. Les réalisations spectrale, solénoïdale, de jauge, cohomologique et déterminantale demeurent des opérations corrélatives de la structure discrète conjointe du continu. L'extraction du registre et les représentations régionales s'effectuent sur ce même domaine. Dans ce qui suit, ce sont toutes des sorties HMT précédemment construites.

\subsection{La collection marquée et son inverse}

Soit \(L\geq2\), et soit \(K=(K_1,\ldots,K_L)\) un registre positif de charge \(Q=\sum_jK_j\). Dans l'application dodécaphasique, \(L=12\) et \(Q=6263\). Nous définissons
\[
 d_j=K_j/Q,\qquad t_0=0,\qquad
 t_j=\sum_{a=1}^{j}d_a,\qquad
 B(K)=\begin{pmatrix}1&\cdots&1\\t_0&\cdots&t_L\end{pmatrix}.
\]
La collection \(\mathcal G_L^{\rm aff}\) est l'image de ces plans de lignes dans \(\operatorname{Gr}_{>0}(2,L+1)\), en conservant les colonnes ordonnées \(0,\ldots,L\), les marques extrêmes et la charge \(Q\). La normalisation affine fixe \(t_0=0,t_L=1\). On note \(\Delta_{ij}\) le mineur des colonnes \(i,j\) de tout représentant du plan.

\begin{proposition}[Inverse de la réalisation marquée]
\label{prop:compat-inversa-marcada}
L'application \(K\mapsto([B(K)],Q)\) est injective. Son inverse sur l'image est donné par
\begin{equation}
 K_j=Q\frac{\Delta_{j-1,j}}{\Delta_{0L}},
 \qquad
 t_j=\frac{\Delta_{0j}}{\Delta_{0L}},
 \quad 1\leq j\leq L.
 \label{eq:compat-inversa-plucker}
\end{equation}
Les quotients sont invariants sous tout changement inversible de base
des deux lignes.
\end{proposition}
\begin{proof}
Dans le représentant déclaré, \(\Delta_{ij}=t_j-t_i\), et \(\Delta_{0L}=1\). Les formules récupèrent les différences consécutives et leurs sommes. Si l'on remplace \(B\) par \(AB\), avec \(A\) inversible, tous les mineurs sont multipliés par \(\det A\) ; leurs quotients restent égaux. La charge récupère l'échelle antérieure à la normalisation. Par conséquent, deux registres ayant les mêmes données marquées ont les mêmes coordonnées.
\end{proof}

La condition d'appartenance à \(\mathcal G_L^{\rm aff}\) importe. Les quotients consécutifs d'un plan positif arbitraire peuvent ne pas satisfaire l'identité \(\sum_j\Delta_{j-1,j}=\Delta_{0L}\). Dans la collection construite, cette identité est télescopique. L'inverse agit sur l'image marquée qui vient d'être définie.

\subsection{Algèbre des lecteurs et changements de coordinatisation}

Soit \(\mathcal D\) l'image d'une représentation de l'état enrichi de la forme \(x\mapsto(K(x),\xi(x))\). Le symbole \(\xi\) réunit exactement les autres coordonnées utilisées par le lecteur : représentations régionales, origine et orientation, route, signature, conditions de frontière et sections dimensionnelles, selon l'objet considéré. On ne suppose pas ces coordonnées indépendantes. On définit
\[
 \mathcal F:\mathcal D\longrightarrow\widetilde{\mathcal G},
 \qquad (K,\xi)\longmapsto([B(K)],Q,\xi),
 \qquad \widetilde{\mathcal G}:=\mathcal F(\mathcal D).
\]
La proposition précédente fournit l’inverse \(\mathcal I\) de
\(\mathcal F\) sur son image.

\begin{theorem}[Représentation fidèle de l'algèbre des lecteurs retenus]
\label{thm:compat-algebra-lectores}
Pour chaque lecteur \(\mathcal R:\mathcal D\to V\), de codomaine
déclaré \(V\), il existe un unique lecteur géométrique
\(\mathcal R^{\rm geom}:\widetilde{\mathcal G}\to V\) tel que
\[
 \mathcal R^{\rm geom}=\mathcal R\circ\mathcal I,
 \qquad \mathcal R^{\rm geom}\circ\mathcal F=\mathcal R.
\]
Pour des lecteurs scalaires, cette correspondance préserve les sommes, les produits, l'unité et la conjugaison. C'est un isomorphisme de l'algèbre des fonctions sur \(\mathcal D\) avec l'algèbre correspondante sur \(\widetilde{\mathcal G}\).
\end{theorem}
\begin{proof}
La composition est définie parce que \(\mathcal I\mathcal F\) est
l’identité. Si deux lecteurs géométriques ont la même composition
avec \(\mathcal F\), ils coïncident sur son image, qui est tout
\(\widetilde{\mathcal G}\). Les opérations algébriques se vérifient
point par point ; par exemple,
\((\mathcal R_1\mathcal R_2)\circ\mathcal I=
(\mathcal R_1\circ\mathcal I)(\mathcal R_2\circ\mathcal I)\).
La composition inverse est la restriction par \(\mathcal F\).
\end{proof}

La fidélité concerne la représentation \((K,\xi)\). Pour un observable \(\mathcal O\) de l'état enrichi complet, le critère de descente est que \(\mathcal O(x)=\mathcal O(x')\) dès que \((K(x),\xi(x))=(K(x'),\xi(x'))\). La nécessité s'obtient par composition ; la suffisance permet de définir la valeur sur une fibre à partir de n'importe lequel de ses représentants. Ce critère conserve explicitement l'information qui doit accompagner le registre.

Soit maintenant \(\mathcal T\) une triangulation du polygone étiqueté déterminant une coordinatisation de Plücker positive. Ses coordonnées de frontière et une normalisation homogène sont incluses, par exemple \(\Delta_{0L}=1\). Si l'on échange la diagonale \(ac\) contre \(bd\), avec \(a<b<c<d\), le changement de coordinatisation est
\begin{equation}
 \Delta_{bd}=
 \frac{\Delta_{ab}\Delta_{cd}+\Delta_{ad}\Delta_{bc}}{\Delta_{ac}}.
 \label{eq:compat-cambio-carta}
\end{equation}
Le dénominateur est positif et la relation de Plücker prouve que les deux coordinatisations évaluent le même point.

\begin{corollary}[Compatibilité des lecteurs avec les mutations]
\label{cor:compat-mutaciones}
Soient \(r_{\mathcal T}\) et \(r_{\mathcal T'}\) les expressions de \(\mathcal R^{\rm geom}\) dans deux coordinatisations conservant les marques, et \(\mu_{\mathcal T\mathcal T'}\) leur changement positif. Alors
\[
 r_{\mathcal T'}\circ\mu_{\mathcal T\mathcal T'}
 =r_{\mathcal T}
\]
sur le domaine commun. Une suite fermée de changements de coordinatisation qui conserve le point marqué conserve tous ces lecteurs.
\end{corollary}
\begin{proof}
Les deux reconstructions au moyen de \eqref{eq:compat-inversa-plucker} récupèrent le même \(K\). Les coordonnées \(\xi\) sont transportées avec leurs marques. Appliquer \(\mathcal R\) prouve l'égalité ; l'itérer prouve l'affirmation concernant une suite fermée.
\end{proof}

Les lecteurs peuvent s'exprimer dans des coordinatisations distinctes tout en restant des fonctions du même objet. Une permutation des étiquettes agit en outre sur la marque positionnelle ; une transformation qui change le point agit sur l'argument du lecteur. Les deux opérations possèdent leur loi de transport spécifique.

L’action diagonale par colonnes fournit un contrôle explicite.
Pour \(\lambda_i>0\),
\(\Delta_{ij}\mapsto\lambda_i\lambda_j\Delta_{ij}\).
Le birapport \(\Delta_{ab}\Delta_{cd}/
(\Delta_{ad}\Delta_{bc})\) est conservé, tandis que
\(\Delta_{j-1,j}/\Delta_{0L}\) reçoit le facteur
\(\lambda_{j-1}\lambda_j/(\lambda_0\lambda_L)\).
Avec \(t=(0,1/4,1/2,1)\) et \(\lambda=(1,2,1,1)\), les
séparations reconstruites seraient \((1/2,1/2,1/2)\), dont la somme
est \(3/2\). Les seuls birapports conservent un quotient différent
de celui exigé par l’inverse marqué.

\subsection{Composition explicite avec l'action et la réponse constitutive}

Dans la coordinatisation dodécaphasique, on écrit \(b=1000\) et
\begin{equation}
 \kappa^{\rm geom}=
 \frac{Q}{b^{12}-1}\sum_{j=1}^{12}
 b^{12-j}\frac{\Delta_{j-1,j}}{\Delta_{0,12}},
 \qquad
 \alpha^{\rm geom}=\pi_{\rm HMT}+e_{\rm HMT}
 -\varphi_{\rm HMT}-4-\kappa^{\rm geom}.
 \label{eq:compat-alpha-geometrica}
\end{equation}
Ce sont les expressions géométriques du lecteur de clôture \eqref{eq:mass-k-alpha}. Les trois représentations régionales et \(K\) partagent une génération coinductive ; l'expression utilise leur ordre constructif interne. Le lecteur d'action \eqref{eq:mass-k-action}, évalué en \(\alpha^{\rm geom}\), détermine \(\eta_{\rm ret}^{\rm geom}\), et donc
\[
 \hbar_{\rm ret}^{\rm geom}
 =\eta_{\rm ret}^{\rm geom}10^{-34}\mathcal U_S
 =\hbar_{\rm ret}.
\]
L’égalité a lieu dans la droite d’action \(\mathcal L_S\)
avec sa section \(\mathcal U_S\) ; elle conserve la décade provenant
du lecteur déterminantal \(\varphi_{\rm HMT}\alpha^{16}\).
La composition ne modifie pas la section et ne l’applique pas une seconde fois.

Dans la coordinatisation angulaire relevée qui conserve l'origine et la branche réelle de la représentation, la paire angulaire ultérieure détermine
\[
 x=\frac{\pi_{\rm HMT}}{180}(1000\alpha),\qquad
 y=\frac{\pi_{\rm HMT}}{90}(\eta_{\rm ret}+\alpha),\qquad
 q_\pm=e^{-x\mp y}.
\]
La branche fait partie des marques retenues : la réduction d'un angle modulo \(360\) à elle seule conserve une autre information. Les égalités qui suivent emploient le relèvement déclaré, dans lequel la première coordonnée est \(1000\alpha\).
Dans la chambre \(x>|y|\), avec des projecteurs de feuillet fixés \(P_\pm\), soit \(T=q_+P_++q_-P_-\). La réponse du corpus compare les incidences \(90\) et \(120\) au moyen de
\[
 \mathcal V=(I-T^{90})(I-T^{120})^{-1}
 =r_+P_++r_-P_-,\qquad
 r_\pm=f(q_\pm^{30}),\quad
 f(s)=\frac{1+s+s^2}{1+s+s^2+s^3}.
\]
La règle vérifie
\(\mathcal V\Sigma_4(T^{30})=\Sigma_3(T^{30})\), où
\(\Sigma_j(S)=\sum_{a=0}^{j-1}S^a\). Son unicité découle
de l’inversibilité de \(\Sigma_4(T^{30})\).
Les quatre lectures normalisées sont
\[
 \widehat\varepsilon=r_+^2,\quad
 \widehat\mu=r_-^2,\quad
 \widehat Z=r_-/r_+,\quad
 \widehat c=(r_+r_-)^{-1}.
\]
Dans les droites dimensionnelles d’action, de charge et de vitesse, elles se réalisent comme
\[
 Z=\widehat Z\,u_Su_Q^{-2},\quad c=\widehat c\,u_C,
 \qquad
 \mu=\widehat\mu\,u_Su_C^{-1}u_Q^{-2},\quad
 \varepsilon=\widehat\varepsilon\,u_S^{-1}u_C^{-1}u_Q^2.
\]
Les identités \(\mu\varepsilon=c^{-2}\) et \(\mu/\varepsilon=Z^2\) s'obtiennent par produit et quotient. Le corollaire des changements de coordinatisation conserve cette réponse complète, avec les marques de feuillet et les sections dimensionnelles incluses dans \(\xi\). Le même argument s'applique à un lecteur de masse dont la route, la signature et l'opérateur de retour sont retenus, sur le domaine de la section~\ref{sec:k-mass-bridge}.

\begin{theorem}[Récupération du registre à partir de la réponse étiquetée]
\label{thm:compat-recuperacion-constitutiva}
Les représentations régionales, l'orientation des feuillets, la base \(b=1000\), l'origine, la longueur douze et le relèvement angulaire déclaré étant fixés, l'application \(K\mapsto(\widehat\varepsilon,\widehat\mu)\) est injective sur les registres entiers de son image physique satisfaisant \(0\leq K_j<b\) et la chambre contractive indiquée. Son inverse exact est la composition des opérations suivantes :
\begin{align*}
 r_+&=\sqrt{\widehat\varepsilon},&
 r_-&=\sqrt{\widehat\mu},\\
 r_\pm s_\pm^3+(r_\pm-1)(1+s_\pm+s_\pm^2)&=0,
 &0<s_\pm<1,\\
 x&=-\frac1{60}\log(s_+s_-),&
 \alpha&=\frac{180x}{1000\pi_{\rm HMT}},\\
 \kappa&=\pi_{\rm HMT}+e_{\rm HMT}-\varphi_{\rm HMT}-4-\alpha,
 &N&=(b^{12}-1)\kappa,\\
 K_j&=\left\lfloor N/b^{12-j}\right\rfloor\bmod b.
\end{align*}
\end{theorem}
\begin{proof}
Les racines positives récupèrent les coefficients étiquetés. La dérivée
\[
 f'(s)=-\frac{s^2(3+2s+s^2)}{(1+s+s^2+s^3)^2}<0
 \qquad(0<s<1)
\]
et ses limites \(1\) et \(3/4\) prouvent que l'équation cubique possède exactement une racine dans cet intervalle. Puisque \(s_\pm=e^{-30x\mp30y}\), leur produit est \(e^{-60x}\). On récupère \(x\), \(\alpha\) et \(\kappa\). Dans l'image considérée, \(N=\sum_jK_jb^{12-j}\) est un entier compris entre zéro et \(b^{12}-1\). La division euclidienne successive récupère de manière unique ses douze chiffres, y compris les zéros initiaux.
\end{proof}

La précision nécessaire à une récupération numérique peut également être exprimée. Avec des représentations régionales exactes, une erreur \(|\widetilde\alpha-\alpha|<1/[2(b^{12}-1)]\) implique \(|\widetilde N-N|<1/2\) ; l'arrondi à l'entier le plus proche restitue \(N\). Si les représentations régionales comportent des erreurs, la même condition s'applique à la somme des erreurs de \(\pi+e-\varphi-\alpha\). L'injectivité exacte et la résolution métrologique sont ainsi des propriétés distinctes.

Le domaine entier est essentiel. Dans un prolongement analytique des registres positifs, la perturbation de trois positions consécutives proportionnelle à \((1,-(b+1),b)\) conserve à la fois la somme et la lecture positionnelle, puisque \(1-(b+1)+b=0\) et \(b^2-b(b+1)+b=0\). Pour des perturbations suffisamment petites, elle conserve la positivité et change le point grassmannien. L'injectivité précédente utilise la condition de douze chiffres entiers ; la collection analytique continue possède d'autres fibres.

\subsection{Raffinement et lecteurs hérités}

Si chaque séparation \(d_j\) est subdivisée en séparations positives \(d_{j,a}\) avec \(\sum_ad_{j,a}=d_j\), les extrémités originales conservent leur position. Soient \(\rho\) l'opération consistant à sommer les successeurs de chaque intervalle et \(H\) l'inclusion de ses extrémités dans le nouvel ensemble ordonné. Alors
\[
 \mathcal I_{\rm heredada}(B(\widetilde d),Q)
 =\bigl(Q\sum_ad_{j,a}\bigr)_j=K,
 \qquad
 \mathcal R_{\rm heredada}=\mathcal R(K,\xi).
\]
L’égalité prouve la commutativité entre reconstruction héritée et
subdivision. Dans une suite de subdivisions, les sommes s’associent ;
le résultat persiste dans chaque composition finie. Cela conserve le
lecteur de longueur douze et ses marques. Un lecteur défini sur les
nouveaux blocs aura son propre domaine et sa loi de comparaison.

\subsection{Formes canoniques à valeurs dans une droite dimensionnelle}

Soit \(P\) l’un des polytopes de rang un construits dans le
manuscrit et \(P=\bigcup_\sigma P_\sigma\) une subdivision
orientée du même support, de multiplicité un et à faces compatibles.
Nous écrivons \(\Omega_\sigma\) et \(\Omega_P\) pour leurs formes
canoniques. Leurs identités sont
\(\Omega_P=\sum_\sigma\Omega_\sigma\) et
\(\operatorname{Res}_F\Omega_P=\Omega_F\), avec la convention
d’orientation fixée dans le développement amplituédrique.
La droite d’action \(\mathcal L_S\) est considérée comme constante sur
ce polytope ; sa section provient de l’état HMT fixé.

\begin{theorem}[Transport de la forme par une section d'action]
\label{thm:compat-forma-accion}
Pour \(a\in\mathcal L_S\), constant par rapport aux coordonnées du polytope,
\[
 \boldsymbol\Omega_P:=a\otimes\Omega_P
 =\sum_\sigma a\otimes\Omega_\sigma,
 \qquad
 \operatorname{Res}_F\boldsymbol\Omega_P=a\otimes\Omega_F.
\]
Les identités se conservent sous des subdivisions successives et des résidus
itérés sur des drapeaux de faces. Elles s’appliquent en particulier à
\(a=\hbar_{\rm ret}\).
\end{theorem}
\begin{proof}
Le produit tensoriel avec une droite est linéaire. Le résidu agit sur le facteur différentiel et commute avec le coefficient constant. La compatibilité des subdivisions et des résidus itérés est transportée par les mêmes opérations linéaires. Dans un changement de base dimensionnelle \(u'_S=\lambda\,u_S\), avec \(\lambda>0\), le coefficient numérique change selon \(\lambda^{-1}\) ; le tenseur et ses résidus restent égaux.
\end{proof}

\begin{theorem}[Condition de recollement des coefficients de frontière]
\label{thm:compat-pegado}
Soient \(a_\sigma\) des sections de la même droite, régulières dans un
voisinage des faces et méromorphes dans une réalisation ambiante commune.
Supposons que leurs produits avec les formes canoniques aient, le
long des faces considérées, des pôles au plus simples. Dans
l’intérieur relatif d’une face commune \(F\) de codimension un de deux cellules
\(\sigma,\tau\), d’orientations opposées,
\begin{equation}
 \operatorname{Res}_F\left(a_\sigma\Omega_\sigma+a_\tau\Omega_\tau\right)
 =\left(a_\sigma|_F-a_\tau|_F\right)\Omega_F.
 \label{eq:compat-defecto-frontera}
\end{equation}
Par conséquent, le pôle de cette paire sur le diviseur de \(F\) disparaît
exactement lorsque les deux traces coïncident. Si les coefficients sont
constants et si le graphe d’adjacence des cellules est connexe, l’annulation
de chaque paire de cellules à travers leur face commune équivaut à l’existence
d’un unique coefficient commun \(a\).
\end{theorem}
\begin{proof}
Avec une coordonnée locale transverse \(z\), une forme logarithmique s'écrit \(dz/z\wedge\omega+\eta\), où \(\eta\) est régulière par rapport à ce diviseur. La régularité de \(a_\sigma\) donne \(\operatorname{Res}_F(a_\sigma\Omega_\sigma)
=a_\sigma|_F\operatorname{Res}_F\Omega_\sigma\). Les deux orientations produisent \(\Omega_F\) et \(-\Omega_F\). La forme \(\Omega_F\) est non nulle dans l'intérieur relatif de la face ; l'égalité des résidus équivaut à l'égalité des traces. Pour un pôle simple, l'annulation du résidu élimine le pôle. Dans le cas constant, l'égalité se propage le long de chaque arête du graphe connexe d'adjacence.
\end{proof}

Pour la somme globale, il faut compter tous les termes ayant un pôle sur le
même diviseur ambiant. Si \(H\) est l’hyperplan qui contient \(F\),
et si \(I_H\) réunit ces termes, l’identité générale est
\[
 \operatorname{Res}_{H}\left(\sum_\nu a_\nu\Omega_\nu\right)
 =\sum_{\nu\in I_H}a_\nu|_H\operatorname{Res}_H\Omega_\nu,
\]
pourvu que les coefficients soient réguliers sur \(H\). Une facette
coplanaire, même si elle n’est pas adjacente à la face réelle \(F\), participe à cette
somme. La formule à deux termes représente le résidu global lorsque
les autres termes sont réguliers sur ce diviseur. L’annulation
par paires de toutes les facettes intérieures élimine leurs contributions ;
les facettes extérieures conservent les résidus de frontière correspondants.

Sur une face de codimension supérieure, cette formule est appliquée successivement avec les orientations induites. Les signes sont ceux de l'incidence cellulaire et satisfont l'annulation du bord du bord. La collection de lecteurs locaux définit un recollement compatible deux à deux lorsque leurs traces coïncident après transport vers la même droite de coefficients. Le résidu global est obtenu en sommant les contributions de chaque diviseur. Cet énoncé fournit un critère effectif pour répartir une lecture physique sur des coordinatisations ou des cellules sans introduire de discontinuités artificielles.

\subsection{Annulation des pôles et normalisation globale}

L'annulation intérieure détermine une condition locale. L'identification à une forme canonique globale requiert également le diviseur des pôles et les résidus extérieurs normalisés. Dans l'espace projectif, deux formes rationnelles de degré maximal ayant les mêmes pôles logarithmiques et les mêmes résidus diffèrent d'une forme holomorphe de degré maximal ; \(H^0(\mathbb P^m,\Omega^m)=0\) démontre leur égalité. Telle est la condition globale utilisée dans la preuve du polytope \cite{positive}.

Un exemple par intervalles montre pourquoi les coefficients variables
exigent cette condition. Pour \(0<c<1\), soient
\[
 \Omega_{0c}=\frac{c\,dx}{x(c-x)},\qquad
 \Omega_{c1}=\frac{(1-c)\,dx}{(x-c)(1-x)},\qquad
 \Omega_{01}=\frac{dx}{x(1-x)}.
\]
Avec \(a_1(x)=1+x(c-x)\) et \(a_2(x)=1\), leurs traces coïncident
à toutes les extrémités pertinentes, et pourtant
\[
 a_1\Omega_{0c}+a_2\Omega_{c1}=\Omega_{01}+c\,dx.
\]
Le terme supplémentaire est régulier dans la coordinatisation affine et possède un pôle à l'infini projectif. Les résidus aux extrémités finies et l'annulation en \(c\) sont conservés, tandis que la forme résultante possède un autre comportement global. Multiplier l'exemple par une section de \(\mathcal L_S\) produit le même contrôle dimensionnel.

La conséquence d'ensemble est précise : la représentation grassmannienne marquée transporte les lecteurs retenus ; les mutations conservent leurs valeurs lorsqu'elles ne changent que la coordinatisation ; le raffinement conserve leurs évaluations héritées ; et les formes à valeurs dans la droite d'action s'assemblent par égalité des traces et contrôle des pôles. Chaque égalité exprime l'opération qui conserve l'objet correspondant.
